\begin{table}[htbp]
\centering
\small
\begin{tabular}{lrrrrrl}
\toprule
Landscape & Distance km & Capacity Direct bps Hz & Capacity HAPS UAV bps Hz & Capacity UAV UE bps Hz & Capacity UAV Relay bps Hz & Bottleneck Hop \\
\midrule
suburban & 0 & 2.67 & 10.20 & 0.00 & 0.00 & UAV-UE \\
suburban & 10 & 2.40 & 10.20 & 0.00 & 0.00 & UAV-UE \\
suburban & 20 & 1.86 & 10.20 & 0.01 & 0.01 & UAV-UE \\
suburban & 30 & 1.37 & 10.20 & 0.09 & 0.09 & UAV-UE \\
suburban & 50 & 0.75 & 10.20 & 0.02 & 0.02 & UAV-UE \\
suburban & 75 & 0.49 & 10.20 & 0.00 & 0.00 & UAV-UE \\
suburban & 100 & 0.46 & 10.20 & 0.00 & 0.00 & UAV-UE \\
high rise & 0 & 2.67 & 10.20 & 0.00 & 0.00 & UAV-UE \\
high rise & 10 & 2.40 & 10.20 & 0.00 & 0.00 & UAV-UE \\
high rise & 20 & 1.86 & 10.20 & 0.01 & 0.01 & UAV-UE \\
high rise & 30 & 1.37 & 10.20 & 0.03 & 0.03 & UAV-UE \\
high rise & 50 & 0.75 & 10.20 & 0.01 & 0.01 & UAV-UE \\
high rise & 75 & 0.49 & 10.20 & 0.00 & 0.00 & UAV-UE \\
high rise & 100 & 0.46 & 10.20 & 0.00 & 0.00 & UAV-UE \\
\bottomrule
\end{tabular}

\caption{Table A.4: Direct (HAPS\textrightarrow UE) vs. relay (HAPS\textrightarrow UAV\textrightarrow UE) link comparison, suburban and high-rise landscapes, as a function of distance from the serving HAPS, azimuth-averaged (best/worst case for the landscape-dependent second hop; full 4-landscape data in \texttt{appendix\_data/UAV\_relay\_path\_comparison.csv}). The relay UAV is static, fixed at the HAPS triangle's centroid (500\,m altitude, 37.5\,km from each HAPS); the 38\,GHz HAPS\textrightarrow UAV feeder hop is therefore a constant, landscape-independent link, while the 2\,GHz UAV\textrightarrow UE access hop uses each user's real, azimuth-dependent distance to that fixed point with the landscape-dependent LoS/NLoS blended model (\S III.C). $C_{\text{relay}} = \min(C_{\text{haps-uav}}, C_{\text{uav-ue}})$, implemented in \texttt{uav\_relay\_two\_hop\_capacity\_bps\_hz()}.}
\label{table:app-uav-relay-comparison}
\end{table}
